\clearpage
\section*{Estensione GUI}

L'estensione aggiunge un'interfaccia desktop realizzata con \textbf{Java Swing}
e \textbf{Java2D}. Consente di apprendere un albero da un file locale oppure da
una tabella MySQL attraverso il server MAP, mostrandone nodi, rami, valori
predetti e statistiche. Il client da console resta disponibile per la
predizione guidata e il caricamento degli archivi.

\subsection*{Architettura e riuso della base}

I sorgenti aggiunti appartengono al package \texttt{estensioni.gui} e non
modificano il codice applicativo della base. Il training locale riutilizza i
package \texttt{data} e \texttt{tree} dell'applicazione locale \texttt{src}, in
\path{project/src/src/}; il training remoto utilizza il server già presente in
\path{project/mapServer/src/}.

I due moduli contengono classi con gli stessi nomi qualificati, ma con modalità
di acquisizione differenti: file nel modulo \texttt{src}, JDBC nel server.
Per questo la GUI include soltanto la base locale nel proprio JAR e
comunica con il server tramite socket, senza mescolare le due implementazioni
nello stesso classpath.

\begin{center}
\begin{tabularx}{\textwidth}{@{}lX@{}}
\toprule
\textbf{Classe} & \textbf{Responsabilità} \\
\midrule
\texttt{RegressionTreeGUI} & Finestra principale, selezione dell'origine,
validazione, avvio del training, stato e dialoghi. \\
\texttt{RegressionTreeAccess} & Adattatore in sola lettura della struttura
privata dell'albero del modulo \texttt{src}. \\
\texttt{VisualTree} & Conversione dell'albero locale e calcolo delle
statistiche della rappresentazione visuale. \\
\texttt{VisualNode} & Titolo, metadati, ramo entrante, figli e coordinate
di un nodo visuale. \\
\texttt{RegressionTreeCanvas} & Disposizione dell'albero, disegno e zoom. \\
\texttt{DatabaseTreeClient} & Sessione TCP e ricostruzione dell'albero remoto
tramite il protocollo della base. \\
\texttt{DatabaseTreeSnapshot} & Risultato dell'ispezione di un nodo remoto:
domanda di split oppure predizione. \\
\texttt{DatabaseTreeResult} & Albero visuale e dettagli testuali restituiti
al termine del training remoto. \\
\bottomrule
\end{tabularx}
\end{center}

\subsection*{Training locale}

\texttt{Data} legge il file \texttt{.dat}; \texttt{RegressionTree} esegue
l'induzione ricorsiva già implementata nella base. La soglia di arresto per
numero di esempi è il 10\% del training set iniziale, calcolato con divisione
intera. In alternativa si genera una foglia quando lo split non produce più
di un ramo; il valore della foglia è la media del target.

\texttt{VisualTree.from()} costruisce una rappresentazione separata, senza
cambiare il modello appreso. La base locale non espone radice e figli:
\texttt{RegressionTreeAccess} legge per riflessione i campi \texttt{root} e
\texttt{childTree}, restituendo una copia dell'array dei figli. Questo conserva
l'API originale, ma lega l'adattatore ai nomi e ai tipi dei campi privati.

\clearpage
\subsection*{Rappresentazione grafica}

Entrambe le origini producono un \texttt{VisualTree}: il pannello di disegno
non deve conoscere né il formato dei file né il protocollo di rete.
\texttt{RegressionTreeCanvas} misura prima la larghezza dei sotto-alberi e
poi dispone ogni padre al centro dello spazio riservato ai propri figli.
Nodi e collegamenti vengono disegnati ricorsivamente con antialiasing.

\begin{itemize}
  \item I nodi interni hanno sfondo nero e mostrano l'attributo di split;
        le foglie hanno sfondo bianco e valore predetto in rosso.
  \item Le etichette dei rami riportano uguaglianze discrete oppure confronti
        continui, conservando gli operatori \texttt{<=} e \texttt{>}.
  \item I valori predetti sono formattati con il punto decimale e fino a
        quattro cifre decimali. Questa formattazione non cambia il modello.
  \item Lo zoom varia dal 50\% al 150\%, in passi di 10 punti percentuali.
        Le barre di scorrimento consentono l'esplorazione; \textbf{Centra}
        riporta la vista verso la radice senza modificare lo zoom.
\end{itemize}

Le statistiche comprendono numero di nodi, foglie e profondità massima,
misurata in archi: una sola foglia ha profondità zero. In modalità locale
sono mostrati anche numero di esempi, attributi e intervalli degli esempi
associati ai nodi. Gli indici si riferiscono ai dati riordinati durante
l'apprendimento, non alle righe originarie del file. La modalità database
non dispone di questi intervalli e usa metadati descrittivi dell'origine.

\subsection*{Training remoto e protocollo}

La GUI si collega al \textbf{server MAP}, non direttamente a MySQL.
\texttt{DatabaseTreeClient} usa un socket TCP con
\texttt{ObjectOutputStream} e \texttt{ObjectInputStream}. I comandi sono
oggetti \texttt{Integer}; nomi e messaggi sono stringhe serializzate.

\begin{center}
\begin{tabularx}{\textwidth}{@{}lXX@{}}
\toprule
\textbf{Codice} & \textbf{Richiesta} & \textbf{Risposta attesa} \\
\midrule
\texttt{0} & Caricamento della tabella; segue il nome. &
\texttt{Table found!}, poi \texttt{OK}. \\
\texttt{1} & Apprendimento e salvataggio sul server. & \texttt{OK}. \\
\texttt{3} & Predizione a partire dalla radice. &
\texttt{QUERY} e domanda; al termine \texttt{QUERY}, \texttt{OK} e valore. \\
\bottomrule
\end{tabularx}
\end{center}

Il protocollo originale non esporta l'intero albero. Per ricostruirlo, la GUI
avvia una predizione per ciascun nodo da ispezionare, ripercorrendo dalla
radice gli indici dei rami che lo identificano. Quando incontra lo split
cercato, ne conserva la domanda e completa comunque la predizione scegliendo
il ramo zero fino a una foglia: in questo modo consuma tutte le risposte e
mantiene sincronizzati gli stream. La visita prosegue ricorsivamente sui
figli. Un nodo foglia viene ricostruito dal valore numerico ricevuto.

Il codice \texttt{2}, usato dalla CLI per caricare archivi, non viene inviato
dalla GUI. Il server continua a salvare \texttt{nomeTabella.dmp} nella propria
directory di lavoro; la GUI non trasferisce tale archivio al client.

\clearpage
\subsection*{Gestione dello stato e degli errori}

La finestra viene creata sull'\emph{Event Dispatch Thread} di Swing.
Un \texttt{CardLayout} alterna i campi per file e database. Le operazioni di
training sono eseguite direttamente negli ascoltatori dei pulsanti; durante
l'elaborazione viene impostato il cursore di attesa. Al successo, la GUI
aggiorna albero e testo, abilita \textbf{Dettagli} e pianifica il centraggio.

Prima della connessione vengono controllati host non vuoto, porta numerica
nell'intervallo 1--65535 e nome di tabella conforme a
\verb|[A-Za-z][A-Za-z0-9_$]*|. Per il file viene verificato che il percorso non
sia vuoto; i controlli sul contenuto sono demandati a \texttt{Data}.

Un errore preliminare mostra un dialogo senza sostituire l'eventuale risultato
precedente. Se invece il caricamento o il training falliscono, l'albero viene
rimosso, i dettagli vengono svuotati e disabilitati e la barra mostra
\texttt{Training non completato}. Il client di rete implementa
\texttt{Closeable} ed è gestito con \emph{try-with-resources}; la chiusura del
socket termina la sessione lato server.

\subsection*{Limiti dell'implementazione corrente}

\begin{itemize}
  \item Training e comunicazione sono sincroni sul thread Swing: con dati
        grandi o rete lenta la finestra può non rispondere temporaneamente.
        Non sono presenti annullamento o avanzamento percentuale.
  \item La ricostruzione remota ripete percorsi dalla radice. Per un albero
        di $N$ nodi e profondità $h$ richiede fino a un ordine di $N h$ scambi
        lungo i percorsi, oltre all'apprendimento. Non è un trasferimento
        diretto del modello.
  \item Il client remoto impone un limite di 10\,000 nodi e timeout di
        30 secondi per connessione e singole letture. Il timeout non è un
        limite alla durata complessiva del training.
  \item La GUI non offre caricamento o salvataggio locale di \texttt{.dmp},
        esportazione del grafico, modifica dei nodi o predizione interattiva
        tramite clic. Per le predizioni guidate resta disponibile la CLI.
  \item I limiti dell'algoritmo base, compresi alcuni casi continui costanti,
        restano invariati. La visualizzazione non aggiunge correzioni o
        parametri di apprendimento.
  \item Il protocollo usa serializzazione Java senza autenticazione o TLS:
        il servizio è destinato a un ambiente didattico fidato, non
        all'esposizione diretta su reti non affidabili.
\end{itemize}

\subsection*{Distribuzione e avvio}

Dalla cartella principale, con Java 8 o successivo installato:

\begin{lstlisting}
java -jar distribution/gui/mapGUI.jar
\end{lstlisting}

Il JAR include la base locale \texttt{src} e le classi della GUI, senza i test.
La classe principale è \texttt{estensioni.gui.RegressionTreeGUI}.
Il JAR non richiede librerie grafiche esterne né il driver
JDBC sul client: Connector/J rimane nella distribuzione del server.
